\chapter{Of Schooling and the Teaching of the Young}
\label{art:schooling}

\articlenote{In which is set down how the silo's children are watched, taught, and tested, from the earliest years of the nursery through the founding recitation and the aptitude that finds each child's place, until the citizen's own majority and entry into labor.}

\section{The Charge of Schooling}

Every child born within the silo is owed instruction sufficient to keep the founding as Article 1 records it, to read this Pact and any writing the silo's work requires, and to come to majority fit for the oath and the labor Article 2 and Article 13 lay upon a citizen. Schooling is not a privilege a household purchases or withholds at its own pleasure; it is a duty the silo owes every child, and a duty every household owes the silo in the child's attending.

\begin{clauses}
\item Every child, from the age this Article sets down until the citizen's sixteenth naming-day, shall be enrolled in schooling; no household may withhold a child from the schooling this Article requires, save where the Head Physician's judgment excuses attendance for illness or for a season.
\item Schooling is provided without charge to any household, its cost borne as a debt of the silo's keeping, the same as the ration under Article 11 and the dwelling under Article 14.
\item The purpose of schooling is fourfold: that every child come to know the founding as Article 1 records it; that every child learn to read and reckon sufficient for a trade and a household; that every child's own aptitude be discovered and weighed, as Article 13 provides; and that every child, upon majority, be fit to swear the oath under Article 2 and bear the labor Article 13 assigns.
\end{clauses}

\section{The Ages of Instruction}

Every observance this section names is reckoned from the child's naming-day, as Article 2 provides, and not from the day of birth.

\begin{clauses}
\item From the third naming-day, a child too young for the schoolroom is kept, while the members of the child's household are at labor, in the nursery under a Nursery Attendant, that the child be watched, fed, and given the earliest teaching of speech, order, and the household's care.
\item From the sixth naming-day --- the age at which a child is judged to have reached basic societal awareness, sufficient to be held, in some measure, to the order and expectation the schoolroom asks of every child alike --- a child is enrolled in the schoolroom proper, and shall attend upon every ordinary working day, save such rest days as this Pact or the child's own household observes.
\item Instruction continues, without interruption save for illness, from the sixth naming-day to the sixteenth, at which naming-day a citizen's schooling ends and the citizen's labor, under Article 13, begins in full.
\item A child who cannot attend the schoolroom by reason of illness, remoteness of floor, or the household's own labor need shall nonetheless be taught, at the least, the founding recitation and the reading it requires before the child's eighth naming-day, by a teacher sent for that purpose or by such other provision as the schooling of that zone arranges.
\end{clauses}

\section{The Founding and the Limits of Instruction}

The founding recitation under Article 1 is the one matter of history this Pact requires every child to be taught, and it is taught whole, exactly as recorded, and no further.

\begin{clauses}
\item Every child shall be taught the founding exactly as Article 1 records it, complete and entire, before the child's eighth naming-day, and no teacher shall add to, subtract from, or embellish that recitation with any further detail of the founders, the manner of the poisoning, or the world before.
\item A teacher who instructs a child in any invented detail of the founding as though it were known fact offends under Article 16, Section 3, and offends gravely, whether the teaching arose from the teacher's own error or from the teacher's deliberate invention.
\item A child's own questions concerning the founding or the time before shall be answered only in the words Article 1 itself provides; where a child's question reaches beyond what Article 1 records, the teacher shall say plainly that the Pact does not provide more, and shall not speculate before the child.
\item This section binds every teacher, every aide, and every other citizen who instructs a child, whether within the schoolroom or without it.
\end{clauses}

\section{The Teachers of the Young}

\begin{clauses}
\item Teaching the young is a skilled trade entered by full shadowing under Article 13, Section 5, and is further an order-adjacent trade under Article 13, Section 6; a citizen shadowing to become a teacher undergoes, beyond the shadowing itself, such vetting as Article 13 there provides, and remains subject to that vetting throughout the citizen's tenure in the trade.
\item A teacher confirmed in the trade is housed, where the teacher's assignment requires it, among the Up Top under Article 14, Section 2, that the teacher be near the schoolroom the teacher serves.
\item A School Aide may be taken on by short shadowing under Article 13, to assist a teacher in the ordinary care and instruction of the youngest children, and to serve the nursery under a Nursery Attendant; an aide does not teach the founding recitation or the course of study under Section 5 without a teacher present.
\item No citizen may instruct a child in the schoolroom, the nursery, or in place of formal schooling, save a citizen confirmed as a teacher or aide under this section, or a parent or guardian teaching a child some matter this Pact does not reserve to the schoolroom, such as a trade skill or a household craft.
\end{clauses}

\section{The Course of Study}

\begin{clauses}
\item Every child shall be taught to read and write, to reckon numbers sufficient for a trade and for the keeping of a household, and to know this Pact's structure and the citizen's own duties under it, that the citizen may swear the oath under Article 2 with genuine understanding.
\item Every child shall be taught the structure of the silo under Article 1, the offices of its governance, and the ordinary safety of the stairwell, the wallscreen, and the floors, that the child grow into a citizen who keeps the silo's order from knowledge and not from fear alone.
\item The course of study is set by Judicial, in consultation with the teachers of every zone, that a citizen taught on one floor be no less prepared than a citizen taught on another; a teacher may add local instruction suited to a floor's own trades, but may not omit any part of the course this section requires.
\item Beyond the tenth naming-day, a child's course of study shall include a survey of the silo's trades, described plainly and without preference, that a child's own aptitude testing under Article 13 be met with some earlier acquaintance with the choices before them.
\item The course of study set under this section does not include instruction in the working of wire, radio, lens, or fine instrument, whatever a child's own aptitude for such work may show; such instruction is given only within Information Technology's own shadowing under Article 13, to a citizen the Department has itself found fit to receive it.
\end{clauses}

\section{Aptitude, the Approach of Majority, and the Passage into Labor}

\begin{clauses}
\item From the twelfth naming-day, the schools under this Article shall assess each child's aptitude for labor as Article 13, Section 2 provides, and shall record the results for the child's own guidance and for the Head of Department's later reference.
\item In a child's final two years of schooling, from the fourteenth naming-day to the sixteenth, a child may be placed, with the household's consent and the relevant Head of Department's leave, in a trial placement of no more than a season's length within a trade of the child's own aptitude or desire, that the child come to know labor's demands before being bound to them. A trial placement is not a shadowing under Article 13, confers no standing in the trade, and may be ended by the Department, the school, or the household at any time.
\item A child's schooling does not end for a trial placement; the child remains subject to this Article's course of study and to the teacher's charge until the citizen's sixteenth naming-day.
\item Upon a citizen's sixteenth naming-day, schooling ends, the citizen swears the oath under Article 2, and the citizen's shadowing or labor assignment begins under Article 13; a trial placement completed under this section may be credited toward the citizen's later shadowing, at the discretion of the Head of the Department in which the trial was made.
\end{clauses}

\section{The Correction and Wellbeing of Children}

\begin{clauses}
\item A child is a ward of the household or floor to which the child is assigned, and holds, under Article 2, Section 1, every duty of a citizen of the child's years, but not the standing to petition Judicial directly; where a child's matter must be brought before Judicial, the household or, in its absence, the teacher shall bring it.
\item Ordinary correction of a child within the schoolroom --- a delay of privilege, additional labor about the schoolroom, a report to the household --- is the teacher's own charge, and is not an offense under Article 17 nor an occasion for Judicial's hearing.
\item A grave breach of the Pact by a child too young for majority is not answered under Article 17 as though the child were a citizen of full standing. Judicial shall hear the matter with the household present, and shall weigh the child's years in fixing whatever correction, restitution, or instruction the matter requires, saving always the gravest offenses named in Article 17, which Judicial may refer upward regardless of the offender's years.
\item A teacher who observes in a child a persistent distress, a defiance of correction beyond the ordinary, or a fixation upon some matter this Pact does not teach, shall report the matter to the Head Physician under Article 19 or to Judicial, that the child's needs be seen to. Such a report is not itself a charge against the child, and creates no mark upon the child's roll under Article 21 save the fact of the report and its resolution.
\end{clauses}

\section{Records of Schooling and Attendance}

\begin{clauses}
\item Every school shall keep a roll of its enrolled children, their attendance, and the results of the aptitude testing under Section 6, and shall report the roll to Judicial not less than once yearly, for entry into the citizen's record under Article 21.
\item A household that fails, without cause recognized under Section 2, to send a child to schooling as this Article requires is answerable under Article 17; repeated failure may be referred to Judicial for a ruling on the household's fitness to keep the child, without prejudice to the household's ordinary standing in all other matters.
\item A citizen's schooling record --- attendance, aptitude testing, and any trial placement under Section 6 --- is kept in the Judicial archive alongside the citizen's other records under Article 21, and is open to the citizen upon majority and to the citizen's own household before it, but to no other citizen save by Judicial's order.
\end{clauses}
